\begin{table}[H]
\centering\small
\caption{A partner-portfolio counterfactual: months of full negative cycle saved by moving five percentage points of the pre-crisis export-share portfolio, priced with the full-cycle coefficients of Table~\ref{tab:main} at each episode's realised partner growth. Negative values are months saved. An accounting exercise with local, partial-equilibrium coefficients, not a general-equilibrium prediction.}
\label{tab:counterfactual}
\begin{tabularx}{\textwidth}{@{}lYYY@{}}
\toprule
Rebalancing & Episodes & Median months & Interquartile range \\ \midrule
Toward own hub, all economies & 219 & $-0.0$ & $[-0.3,\ 0.1]$ \\
\quad US-hub economies & 69 & $-0.5$ & $[-0.6,\ -0.3]$ \\
\quad EU-hub economies & 146 & $0.0$ & $[-0.1,\ 0.1]$ \\
\midrule
From China to own hub (China share $>0.5$\%) & 167 & $-0.4$ & $[-0.7,\ -0.2]$ \\
\quad US-hub economies & 60 & $-0.5$ & $[-0.8,\ -0.3]$ \\
\quad EU-hub economies & 107 & $-0.4$ & $[-0.7,\ -0.2]$ \\
\bottomrule
\end{tabularx}
\\[2pt]\footnotesize{The move toward the hub is funded proportionally from all non-hub partners (upper panel) or from the China share alone, capped at that share (lower panel). The 219 observations in the upper panel correspond to the 219 unique country-by-peak-year export portfolios across the 263 episodes (69 US-hub, 146 EU-hub, and 4 China-hub in Mongolia); the lower panel restricts to the 167 portfolios with a pre-downturn China export share exceeding 0.5 percent (excluding China-hub episodes).}
\end{table}
